% !TEX root = ../notes.tex

\documentclass[../notes.tex]{subfiles}

\begin{document}

\section{April 15}
Here we go.

\subsection{Reductive and Linear Groups}
Intuitively, a reductive group should be a semisimple group times a torus. We take this as a definition.
\begin{definition}[reductive]
	A connected complex Lie group $G$ is \textit{reductive} if and only if it takes the form
	\[\left((\CC^\times)^r\times G_{\mathrm{ss}}\right)/Z,\]
	where $G_{\mathrm{ss}}$ is a semisimple complex Lie group, and $Z$ is some finite central subgroup.
\end{definition}
\begin{remark}
	Note that $G$ being reductive implies that $\mf g$ is reductive, but the converse does not hold. Indeed, one can show that the Lie group $G=\CC$ is not reductive, which has abelian and hence reductive Lie algebra.
\end{remark}
These reductive groups have maximal compact subgroups.
\begin{notation}
	Let $G=\left((\CC^\times)^r\times G_{\mathrm{ss}}\right)/Z$ be a connected complex Lie group. Then we define
	\[G^c\coloneqq\left(\left(S^1\right)^r\times G_{\mathrm{ss}}^c\right)/Z.\]
\end{notation}
\begin{remark}
	One can show that $G^c$ is independent of the presentation of $G$, simply by perturbing this presentation accordingly.
\end{remark}
\begin{remark} \label{rem:restrict-reps}
	The restriction functor from holomorphic finite-dimensional representations of $G$ to continuous finite-dimensional representations of $G^c$ is an equivalence. Indeed, one shows this because a holomorphic representation splits as a product of a representation of the torus and a representation of the semisimple group, from which one can use \Cref{cor:semisimple-restriction} to handle the semisimple part. For the torus, one needs to know that holomorphic representations of $\CC^\times$ diagonalize as a direct sum of characters which look like $n\colon\CC^\times\to\CC^\times$.
\end{remark}
\begin{remark}
	Choose complex reductive groups $G_1$ and $G_2$. From \Cref{rem:restrict-reps}, it turns out that every homomorphism $G_1^c\to G_2^c$ extends uniquely to a holomorphic homomorphism $G_1\to G_2$. It turns out that the category of connected compact Lie groups embeds faithfully into the category o connected complex reductive groups, whose inverse on objects is $G\mapsto G^c$.
\end{remark}
While we're here, we say something about linear groups.
\begin{definition}[linear]
	A connected Lie group $G$ is \textit{linear} if and only if $G$ admits a faithful finite-dimensional representation. Equivalently, $G$ is linear if and only if it can be realized as a Lie subgroup of $\op{GL}_n(\RR)$ for some $n$.
\end{definition}
\begin{example}
	Any compact group admits a faithful representation by \Cref{prop:compact-has-faithful-rep}, so they are linear.
\end{example}
\begin{example}
	Any complex reductive group is linear. We don't have to worry about the torus (which can be embedded faithfully separately), and then one can pass to the compact part by \Cref{cor:semisimple-restriction}, so we are done by the previous example.
\end{example}
\begin{nex}
	Real semisimple groups need not be linear. For example, the universal cover of $\op{SL}_2(\RR)$ is not linear, by the argument of \Cref{rem:sp-tilde-not-linear}.
\end{nex}
More generally, we have the following.
\begin{proposition}
	Let $\mf g_\theta$ be a real form of a complex semisimple Lie algebra $\mf g$, and let $G$ be a connected complex Lie group with Lie algebra $\mf g$, and set $G_\theta\coloneqq G^{\omega_\theta}$. Then $G_\theta^\circ$ is linear, and every semisimple linear group is of this form.
\end{proposition}
\begin{proof}
	We already know that $G_\theta^\circ$ is linear because it lands in the semisimple complex group $G$, which is then reductive. For the converse, we basically repeat the proof of \Cref{rem:sp-tilde-not-linear}. The idea is that any faithful representation of $G_\theta^\circ$ goes down to the Lie algebra and then lifts to a representation of $G^c$.
\end{proof}

\subsection{Maximal Tori}
Here is our definition.
\begin{definition}[Cartan]
	A \textit{Cartan subalgebra} $\mf h^c$ of $\mf g^c$ is a maximal abelian subalgebra of $\mf g^c$. We let $H^c$ be the exponential of $\mf h^c$ in $G^c$.
\end{definition}
\begin{remark}
	We claim that $\mf h^c\subseteq\mf g^c$ is a Cartan subalgebra if and only if $\mf h^c_\CC$ is a Cartan subalgebra of $\mf g$. Indeed, all elements of $\mf g^c$ are automatically semisimple, so $\mf h^c_\CC$ succeeds at being a torus in $\mf g$, and maximality transfers between the $\mf g^c$ and $\mf g=\mf g^c_\CC$ because maximal just means that the dimension agrees with the rank.
\end{remark}
\begin{remark} \label{rem:exp-to-torus}
	Note that $H^c$ is certainly a torus in $G^c$, and it is also maximal. In fact, $H^c$ is maximal abelian subgroup: any such commuting element $g$ needs to preserve $H^c$ while also preserving the roots, and in fact, $g$ must act by a scalar in $\op U(1)$ on the root spaces. One can then undo the exponential as we did in \Cref{prop:compute-weyl-group} to show that $g\in H^c$.
\end{remark}
\begin{remark}
	As a sanity check, we note that $H\coloneqq\exp(\mf h^c_\CC)$ has $H\cap G^c=H^c$. Indeed, $H\cap G^c$ is some abelian subgroup containing $H^c$.
\end{remark}
All these Cartan subalgebras are conjugate.
\begin{lemma} \label{lem:compact-cartan-conj}
	Any two Cartan subalgebras in $\mf g^c$, equipped with a set of positive roots, are conjugate by $G_{\mathrm{ad}}^c$.
\end{lemma}
\begin{proof}
	Choose two such pairs $(\mf h,\Pi)$ and $(\mf h',\Pi')$. Certainly there is some $g\in G_{\mathrm{ad}}$ with $\op{Ad}_g(h,\Pi)=(h',\Pi')$, so it only remains to descend to $G_{\mathrm{ad}}^c$, for which we use the polar decomposition. On one hand, $\omega(g)=\overline g$ also satisfies $\op{Ad}_{\overline g}(\mf h,\Pi)=(\mf h',\Pi')$, so $\overline g^{-1}g$ commutes with $\mf h^c$ while fixing $\Pi$. Thus, $\overline g^{-1}g$ must live in $H$, so $\overline gh=g$ for some $h\in H$.

	On the other hand, \Cref{thm:polar-decomp} lets us write $g=kp$ for $k\in G^c$ and $p\in P$, and $\overline g=kp^{-1}$ due to the definition of $K$ and $P$. It follows that
	\[kp^{-1}h=kp,\]
	which is then equivalent to $h=p^2$. Thus, $p$ has positive eigenvalues, so $h$ does as well, so we can take a square root of $h$ (in $P$). Thus, $p$ actually commutes with $\mf h^c$ and preserves $\Pi$. We conclude that actually $\op{Ad}_k(\mf h,\Pi)=(\mf h',\Pi')$.
\end{proof}
\begin{remark}
	This conjugacy property is special for the compact. For example, the split form $\mf{sl}_2(\RR)$ has two different tori, given by $\mf{so}_2(\RR)$ and the diagonal $\mf{gl}_1(\RR)$.
\end{remark}
Here is a nice structural result.
\begin{theorem} \label{thm:find-element-in-torus}
	Fix a connected compact Lie group $K$. Then every element is contained in a maximal torus, and all maximal tori are conjugate when equipped with a set of positive roots.
\end{theorem}
\begin{proof}
	By \Cref{prop:classify-compact}, we may assume that $K$ is semisimple, so $K$ takes the form $G^c$ for some semisimple $G$. The second statement now follows from \Cref{lem:compact-cartan-conj} once we realize that $\mf h^c\mapsto H^c$ is a bijection of Cartan subalgebras with maximal tori. Conversely, any maximal torus $T\subseteq K$ induces a maximal abelian subalgebra $\op{Lie}T\subseteq\op{Lie}K$, and the exponential of $\op{Lie}T$ in $K$ must recover $T$ by \Cref{rem:exp-to-torus}.

	It remains to show the first statement. Choose a maximal torus $T$, and let $\varphi\colon K\times T\to K$ be the action map
	\[\varphi(k,t)\coloneqq ktk^{-1}.\]
	Notably, the image of $\varphi$ consists of those elements of $K$ which can be conjugated into $T$, which the above paragraph explains is precisely those elements which are contained in a maximal torus. Notably, the source of $\varphi$ is compact, so the image is closed, so $K\setminus\im\varphi$ is open.

	It is thus enough to show that a dense subset of $K$ can be found in a maximal torus. To this end, for $x\in K$, let $P_x(z-1)$ be the characteristic polynomial of $\op{Ad}_x$ acting on $\mf k$. Then the lowest power of $z$ is the dimension of the centralizer $\mf z_x\subseteq\mf g$. Now, let $a(x)$ be the coefficient of $z^{\op{rank}\mf g}$ in $P_x(z)$. Note that this function is a polynomial in the coordinates of $\op{GL}(\mf g)$, but $a$ vanishes on $K\setminus\im\varphi$, so either $K\setminus\im\varphi$ is empty or $a=0$.
	
	Thus, it remains to show that $a$ is nonzero. Well, for generic $X\in\mf g^c$, one knows that $X$ is regular, meaning that its centralizer $\mf z_X$ has dimension $\op{rank}\mf g^c=\op{rank}\mf g$; this means that $\mf z_X$ is a Cartan subalgebra. Thus, $x=\exp(X)$ commutes with the maximal torus $\exp(\mf z_X)$, so $x\in\exp(\mf z_X)$.
\end{proof}
\begin{corollary}
	Fix a connected compact Lie group $K$ with Lie algebra $\mf k$. Then the exponential map $\exp\colon\mf k\to K$ is surjective.
\end{corollary}
\begin{proof}
	This follows from \Cref{thm:find-element-in-torus} because our maximal tori were the image of an exponential by construction.
\end{proof}
\begin{remark}
	This is not true for complex reductive groups. For example, it is not true for $\op{SL}_2(\CC)$, which we will show on the homework. As an aside, we note that the exponential is surjective for $\op{GL}_n(\CC)$, which one can show using the Jordan normal form.
\end{remark}
While we're here, let's say something about semisimple and unipotent elements.
\begin{defihelper}[semisimple, unipotent] \nirindex{semisimple} \nirindex{unipotent}
	Fix a connected complex reductive group $G$. Then $g\in G$ is \textit{semisimple} (respectively, \textit{unipotent}) if and only if the action of $g$ on any representation is given by a semisimple (respectively, \textit{unipotent}) operator.
\end{defihelper}
\begin{remark}
	On the homework, we will show that it is enough to act by a semisimple (respectively, unipotent) operator on a chosen faithful representation.
\end{remark}
\begin{remark}
	On the homework, we will show that there is a Jordan decomposition: any element $g\in G$ can be written uniquely as a product of a semisimple and a unipotent element.
\end{remark}
\begin{remark}
	It turns out that the exponential defines a homeomorphism between the nilpotent cone of $\mf g$ and the unipotent subvariety of $G$.
\end{remark}

\subsection{The Cartan Decomposition}
We need to start by classifying maximal abelian subspaces of $\mf p$.
\begin{proposition}
	Fix a real form $\mf g_\theta$ of a semisimple complex Lie algebra $\mf g$, and let $\mf g=\mf k\oplus\mf p$ be the usual decomposition into eigenspaces of $\theta$, and similarly $\mf g^c=\mf k^c\oplus\mf p^c$ and $\mf g_\theta=\mf k^c\oplus\mf p_\theta$. Lastly, let $K^c$ be the Lie group of $\mf k$ in $G^c$.
	\begin{listalph}
		\item Let $\mf a$ be a maximal abelian subspace of $\mf p_\theta$. Then the centralizer of $\mf z$ of $\mf a$ in $\mf g^c$ takes the form $\mf m\oplus\mf a$, where $\mf m\subseteq\mf k^c$ is reductive. In fact, if $\mf t\subseteq\mf m$ is a Cartan subalgebra, then $\mf t\oplus\mf a$ is a Cartan subalgebra of $\mf g_\theta$, and $\mf t\oplus i\mf a$ is a Cartan subalgebra of $\mf g^c$.
		\item For generic $a\in\mf a$, the centralizer of $a$ in $\mf p_\theta$ is $\mf a$.
		\item For all $p\in\mf p_\theta$, there is $k\in K^c$ for which $\op{Ad}_kp\in\mf a$.
		\item All maximal abelian subspace of $\mf p_\theta$ are conjugate by $K^c$.
	\end{listalph}
\end{proposition}
\begin{proof}
	Here we go.
	\begin{listalph}
		\item Choose $X\in\mf z$. Then we may write $X=X_+\oplus X_-$ where $X_+\in\mf k^c$ and $X_-\in\mf p_\theta$. Now, $[X_+,\mf a]=0$ automatically, so it follows that $[X_-,\mf a]=0$ as well, so $X_-\in\mf a$ because $\mf a$ is maximal. Thus, the collection of $X_+$ coordinates allow us to decompose $\mf z$ into $\mf m\oplus\mf a$ for some Lie subalgebra $\mf m\subseteq\mf k^c$. Note that $\mf m$ is reductive because the Killing form on $\mf k$ restricts to a non-degenerate invariant bilinear form.

		Lastly, choosing a Cartan subalgebra $\mf t\subseteq\mf m$, we see that $\mf t\oplus i\mf a$ is a maximal subalgebra of in $\mf g^c$: indeed, anyone $X_++X_-$ commuting with $\mf t\oplus i\mf a$ has $X_-\in\mf a$ as above, so $X_+\in\mf m$, but then $[X_+,\mf t]=0$ implies $X_+\in\mf t$ by its maximality. The argument that $\mf t\oplus\mf a$ is a Cartan subalgebra of $\mf g_\theta$ is similar.

		\item Write $T_{\mf a}\coloneqq\exp(i\mf a)$, which is contained in $G^c$. Then $T_{\mf a}$ is some compact torus, which looks like some power of $S^1$. For generic $a$, the one-parameter subgroup in $T_{\mf a}$ is dense. Now, for $p\in\mf p_\theta$ which commutes with $a$, we conclude that $p$ commutes with $T_{\mf a}$. Taking the differential implies that $p$ commutes with $\mf a$.

		\item This is the hardest part. Choose $a\in\mf a$ generic as in (b). Thus, $\op{Ad}_k(p)\in\mf a$ is equivalent to $[\op{Ad}_k(p),a]=0$. The trick is to consider the function $f\colon K^c\to\RR$ defined by
		\[f(b)\coloneqq(\op{Ad}_bp,a).\]
		This function is differentiable (because the inner product can be defined by the Killing form, which is algebraic), so we may select $k$ so that $f$ is maximized at $k\in K^c$; set $p'\coloneqq\op{Ad}_kp$, which we claim will do the trick. Well, the derivative $df$ vanishes at $k$. On the other hand, for $x\in\mf k^c$, this derivative $df_k(x)$ is $([x,p'],a)$, so we see that
		\[(x,[p',a])=0\]
		for all $x$, so $[p',a]=0$ by non-degeneracy, so we are done!
	\end{listalph}
\end{proof}

\end{document}